\documentclass[11pt,a4paper]{article}
\usepackage[review]{acl}
\usepackage{amsmath,amssymb,booktabs,array,multirow}
\usepackage{graphicx}
\usepackage{xcolor}
\usepackage{microtype}
\usepackage{hyperref}
\usepackage{url}

% FORMAT NOTE: Built with the current official ACL review template for ARR.
% Recheck the NAACL 2027 venue-specific CFP before submission.

\title{Reachability-Targeted External-State Initialization for On-Policy Distillation:\\
A Falsifiable Protocol and Audited Agent Case Study}

\author{Anonymous submission}

\newcommand{\opd}{\textsc{opd}}
\newcommand{\core}{\textsc{Core-3}}
\newcommand{\rkl}{\mathrm{rKL}}
\newcommand{\todoresult}[1]{\textcolor{red}{#1}}

\begin{document}
\maketitle

\begin{abstract}
On-policy distillation supervises states visited by the student, but long-horizon
agents may rarely reach external world states where a teacher has a measurable
conditional advantage. We study a 2B-parameter language model
acting for six hours through 17 typed tools in an open-source MMORPG. In one
reported run per arm, natural-visitation distillation changes tool use and
interaction tempo but leaves the quest score at 12/30. A second round augments
student rollouts with hand-selected, verified intermediate world states. Evaluated from a fresh,
unseeded world, the resulting policy scores 15/30; all three clustered prompt
variants pass a wall that none passes under the reported base or first-round
policy. This is a weights-only observation, not a replicated causal estimate:
the training rounds differ in more than state source. A targeted teacher-forcing
probe also finds that malformed history can locally reverse the teacher's syntax
preference. In rewritten runtime logs, correction markers do not repeat within a
session. These observations motivate a confirmatory six-arm comparison in which
training-only external states are selected by frozen student-visitation, teacher-
advantage, validity, and recoverability criteria. Random-valid, progress-matched,
teacher-success-prefix, natural-OPD, and guided-OPD controls separate targeted
coverage from generic resets and prefix curricula. We present a falsifiable
protocol and audited case study, not an established method effect.
\end{abstract}

\section{Introduction}

Post-training a small language model to act over long horizons is not just a
token-prediction problem. The model chooses which environment states become
training data. If it never reaches a state where its teacher knows what to do,
dense teacher feedback is unavailable exactly where it is needed. This creates
a state-visitation bottleneck: improving token-level agreement on common states
can alter style without transferring the competence that distinguishes the
teacher.

We examine that bottleneck in a persistent tool-use setting. A Qwen-family 2B
student plays a tile-based MMORPG through a typed action interface. The world
persists across conversation rollovers, and benchmark progress requires long
chains of gathering, navigation, dialogue, and quest actions. A scaffolded 4B
teacher reaches states that the 2B student does not. We apply reverse-KL
on-policy distillation to the student's own emitted tokens. The historical
comparison contrasts natural student visitation with a mixture that begins some
rollouts from verified intermediate database states; the proposed confirmatory
method additionally freezes a measurable state-selection rule before training.

The present evidence is hypothesis-generating and under-replicated. It contains
one complete six-hour run per principal arm; the three agents within a run are
prompt variants sharing weights, harness, and launch conditions, not independent
experimental units. We therefore separate observations from claims throughout.
The observed sequence is: base 12/30, natural-visitation round 12/30, and
state-augmented round 15/30. In the last arm, every prompt variant passes the
same wall from an unseeded evaluation world. This supports the hypothesis that
training visitation matters, but it does not isolate state seeding from all
round-to-round changes or estimate population-level uncertainty.

This working draft makes three bounded contributions:

\begin{enumerate}
  \item We define \emph{reachability-targeted external-state initialization}:
  training-only resets into legally reachable persistent states selected for
  low natural student occupancy, teacher conditional advantage, and
  recoverability. We pair it with controls that can falsify the method claim.
  \item We distinguish behavioral imitation from competence: the first round
  closely changes measured tool mix and tempo toward the teacher while the
  benchmark score remains fixed.
  \item We expose an interface failure that token-average analyses obscure: in
  a targeted teacher-forcing probe, malformed history reverses the teacher's
  local preference. We separate this grading-side observation from a runtime
  recovery result and specify the experiments required to test generality.
\end{enumerate}

The broader lesson is methodological. Model weights, visited states, and
runtime affordances are different interventions. A credible agent-training
result must preserve and cross them rather than roll them into one headline.

\section{Setting}

\subsection{Persistent tool-use benchmark}

The environment is Kaetram, an open-source tile-based MMORPG backed by a game
server and MongoDB. Each policy controls three agents with different priority
prompts (grinder, completionist, and explorer) through 17 typed tools. Tools
cover observation, navigation, combat, gathering, dialogue, quest inspection,
inventory use, and crafting. The prompt includes procedural quest guidance, so
the benchmark measures execution of a supplied plan rather than autonomous
world-model discovery.

The primary development metric, \core{}, counts newly reached stages in three
quest chains: Foresting (3 stages), Herbalist's Desperation (3), and Rick's Roll
(4). Each agent can gain at most ten stages; a three-agent run therefore scores
in $[0,30]$. Runs begin from a freshly initialized world and continue for six
hours. Conversation sessions roll over under a fixed context budget while game
state persists. The intended confirmatory experimental unit is a full run from
a fresh world. Prompt variants are clustered observations within that run.

\subsection{Agent and interface contract}

The policy acts in an observe--reason--act loop through an OpenAI-compatible
model endpoint and a typed tool server. Navigation and selected interactions
contain deterministic affordances such as path finding and automatic walking;
the model still decides goals, tools, and arguments. The current development
record contains an important interface asymmetry: older base-model runs used a
native model-visible tool schema while some fine-tuned runs did not receive an
identical serialization. We therefore do not use those comparisons as primary
evidence. The confirmatory implementation requires normalized full-schema hash
parity across training and every evaluation arm.

\section{On-Policy Distillation Under Visitation Constraints}

Let $\pi_S$ and $\pi_T$ denote student and teacher policies, and let
$d^{\pi_S}$ be the serving-state distribution induced by student rollouts. A
reverse-KL on-policy objective can be written
\begin{equation}
\begin{aligned}
\mathcal{L}_{\mathrm{OPD}}
&=\mathbb{E}_{s\sim d^{\pi_S}}
  \mathbb{E}_{a\sim\pi_S(\cdot\mid s)} \\
&\quad\left[\log\pi_S(a\mid s)-\log\pi_T(a\mid s)\right].
\end{aligned}
\label{eq:opd}
\end{equation}
The teacher scores the student's emitted tokens. Training uses a clipped
importance-weighted update around the rollout policy, with advantages frozen at
data-build time. Full implementation details appear in Appendix~\ref{app:training}.

Equation~\ref{eq:opd} makes the bottleneck explicit. If a teacher advantage is
concentrated on a set of states $H$ for which
$\Pr_{s\sim d^{\pi_S}}[s\in H]\approx0$, token-dense grading supplies almost no
information about that advantage. We change the state source, not the loss:
some student rollouts begin from verified intermediate environment states, and
their records are mixed with natural student rollouts. The student still acts
and the same teacher still grades every emitted token.

The seeded states are typed game states rather than textual demonstrations.
Player rows specify quest stage, skill level, inventory, and a verified walkable
position. End-to-end checks confirm that the game loads each state and exposes
the intended condition through the same tools used at evaluation. Evaluation is
always unseeded. This resembles reset-based curricula in reinforcement learning
\citep{salimans2018montezuma,resnick2018backplay,ecoffet2021goexplore}. TCOD
uses teacher-success prefixes to initialize multi-turn OPD curricula
\citep{wang2026tcod}; Guided-OPD mixes teacher and student turns and anneals
guidance to zero \citep{li2026guidedopd}; ReOPD treats prefix selection as a
teacher-reliability problem \citep{liao2026reopd}. These works preempt generic
novelty for intermediate starts or state-distribution design. Our narrower
object is a direct, prefix-independent external-state initializer. A candidate
state need not lie on a successful teacher trajectory and is selected by a
frozen rule rather than by progress alone.

Let $v(s)$ be estimated natural student visitation, and let $p_T(s)$ and
$p_S(s)$ be repeated teacher and student success rates when initialized at
$s$. The confirmatory selector admits a state only if it passes world invariants,
is unfinished and task-relevant, is recoverable, satisfies $v(s)\leq\tau_v$,
and satisfies $p_T(s)-p_S(s)\geq\tau_\Delta$ together with a minimum teacher
success threshold. Counts and denominators, source-run identifiers, complete
snapshots, and hash-verified validation artifacts are frozen before outcomes.
This rule distinguishes the proposed method from the hand-selected milestone
used in the exploratory round.

\section{Evidence Protocol}

We grade evidence before interpreting it. A \emph{run} is the independent unit;
agents within a run share the policy, launch, infrastructure, and analysis
pipeline. Exact scores from single runs are descriptive. Counts over agents or
sessions diagnose mechanisms but cannot be treated as independent replications
of a policy effect.

The principal comparison uses the same 2B base checkpoint, six-hour unseeded
evaluation protocol, benchmark, and nominal harness. Round 1 trains on natural
base-policy rollouts. Round 2 initializes from round 1 and mixes natural
round-1 rollouts with 2,326 records collected by the round-1 policy from a
verified Herbalist milestone state. The round-2 corpus contains 7,849 states,
split into 7,024 training and 825 held-out records. Because round 2 also changes
the initialization and data batch, round 1 versus round 2 is not a clean
seeded-versus-unseeded training ablation.

\begin{table*}[t]
\caption{Observed development runs and the claims they permit. Scores are
descriptive because there is one run per arm. Recovery is a runtime interface
affordance.}
\label{tab:arms}
\centering
\small
\setlength{\tabcolsep}{5pt}
\begin{tabular}{lccccl}
\toprule
Policy & Training states & Recovery & \core{} & Wall & Interpretation \\
\midrule
Base 2B & -- & off & 12 & 0/3 & Reference run \\
OPD round 1 & natural & off & 12 & 0/3 & Behavior changed; score did not \\
OPD round 2 & natural + seeded & off & 15 & 3/3 & Weights-only, unseeded observation \\
Round-2 weights & -- & on & 17 & -- & Recovery ablation; one run \\
Round-3 weights & broader seeded & on & 18 & -- & Weights plus interface; not pure weights \\
\bottomrule
\end{tabular}
\end{table*}

The \emph{wall} column measures passage of Herbalist stage 1 by the three prompt
variants. It is a useful within-run consistency check, not a binomial sample of
independent policies. We preserve raw pre-rewrite emissions for confirmatory
format analyses; rewritten history is not an authoritative defect counter.

\section{Exploratory Results}

\subsection{Natural-visitation OPD transfers behavior, not score}

Round 1 leaves \core{} unchanged at 12/30 and every agent reaches the same 4/10
stage profile as base. Nevertheless, several behavioral measurements move
toward the 4B teacher (Table~\ref{tab:behavior}). The average number of turns per
session exactly matches the teacher, observation and navigation shares approach
the teacher, and query frequency overshoots it. The result is consistent with
mode-seeking behavioral imitation, while showing that agreement on frequently
visited states is insufficient for this quest wall.

\begin{table}[t]
\caption{Behavior changes in the first OPD round. Percentages are shares of
tool calls in one run per policy; they are diagnostics, not independent trials.}
\label{tab:behavior}
\centering
\small
\begin{tabular}{lrrr}
\toprule
Measure & Base & Round 1 & Teacher \\
\midrule
Turns per session & 9.1 & \textbf{6.1} & 6.1 \\
Observe share & 26.6\% & \textbf{35.5\%} & 36.1\% \\
Navigate share & 14.3\% & \textbf{9.4\%} & 9.1\% \\
Attack share & 17.9\% & \textbf{6.6\%} & 9.1\% \\
Query-quest share & 5.3\% & \textbf{18.0\%} & 8.7\% \\
Chars./turn & 411 & \textbf{1,248} & 864 \\
\bottomrule
\end{tabular}
\end{table}

Per-verb reverse-KL signs and behavioral movement are also aligned in the
development logs: verbs whose held-out divergence falls generally move toward
the teacher's usage, while query-quest both diverges and behaviorally overshoots
the teacher. This is descriptive triangulation. It does not establish that a
specific token update caused downstream quest behavior.

\subsection{The seeded-training round is followed by unseeded wall passage}

The 4B teacher demonstrates competence beyond the Herbalist stage-1 wall, while
the base 2B almost never visits the required state naturally. Seeded collection
starts the round-1 policy at the accepted quest with the required gathering
skill and near the target resource. All three seeded collectors pass the wall,
creating student-authored actions and teacher scores in this region.

After training on the natural-plus-seeded mixture, round 2 is evaluated from a
fresh unseeded world with recovery disabled. It scores 15/30, with each prompt
variant at 5/10. All three pass Herbalist stage 1, compared with zero of three
under both the base and round-1 policies. Logs show earlier quest acceptance,
explicit tracking of the skill gate, resource rotation after empty harvests,
and less premature return to the quest giver. These are concrete trajectory
differences consistent with the intended mechanism.

The valid statement is deliberately narrow: \emph{a seeded-training round was
followed by unanimous unseeded wall passage in one run}. The present data do not
show that seeding alone caused the improvement. The required causal experiment
trains fresh students from the same checkpoint under identical budgets, seeds,
and data construction, varying only whether the verified state mixture is
included (Section~\ref{sec:confirmatory}).

\subsection{Weights and recovery solve different failures}

Round 3 combines broader state seeding with a runtime recovery affordance and
scores 18/30. That number is not a pure-weights result. Applying the same
recovery affordance to round-2 weights produces 17/30 in one run, suggesting
that most of the round-2-to-round-3 stage difference may come from the
interface. The clean weights-only comparison remains base 12/30 versus round 2
15/30, one run per arm.

The recovery mechanism detects a dropped malformed call, reconstructs the
executable call, rewrites the retained history into canonical form, executes the
action, and returns an explicit correction note. In the audited round, 405
sessions contain a recovery marker; each contains exactly one such marker and
no logged recurrence after the correction. This is a system-counting fact. It
does not show that weights learned the format or that recovery generalizes to
other defects.

\section{A Teacher-Forcing Failure at the Interface}

Round 1 introduces a malformed parameter-key pattern that is absent in the
base run. The emitted form places a Boolean value inside the parameter key, so
the parser silently drops the intended argument. The training records contain
the correct syntax, and the teacher favors the correct delimiter in clean
base-policy contexts. A targeted two-sided probe, however, gives a different
result after the malformed prefix is already present: at the divergence token,
the 4B teacher assigns log probability $-0.161$ (about 85\% probability) to
continuing the malformed form, versus $-0.253$ under the student. On a matched
correct-form state, the preference favors the correct continuation.

Thus teacher generation and teacher-forced grading are distinct objects. A
teacher that generates canonical syntax from a clean history can locally reward
continuation of a student's defect after that defect enters the context. Under
dense on-policy distillation, this creates wrong-signed feedback at precisely
the visited malformed prefix. We call this a \emph{teacher-forcing copy-prior
hypothesis}. The current evidence is one targeted defect family and a small
probe, so the name describes a candidate mechanism rather than a general law.

This observation relates to exposure bias and to failures of token-level
teacher agreement on degraded prefixes. It also gives a testable prediction:
paired histories differing only in canonical versus malformed prior syntax
should change teacher log-odds while holding state and candidate continuations
fixed. Section~\ref{sec:confirmatory} turns that prediction into a multi-state,
multi-defect experiment.

\section{Confirmatory Study Design}
\label{sec:confirmatory}

The repository contains launch and validation work for the following design;
results are not reported until immutable run bundles exist.

\paragraph{Primary endpoint and unit.}
The independent unit is a six-hour run from a fresh seeded world. The
preregistered primary endpoint is total \core{} stage gain across the three
clustered agents. Quest-wall passage is secondary. All runs, confidence
intervals, effect sizes, and stopping rules must be reported.

\paragraph{State-source ablation.}
Train fresh students from the same checkpoint under five matched initial-state
distributions: natural OPD; targeted states; random valid states; progress-
matched states; and TCOD-B2F teacher-success-prefix states. Add visitation-only
and teacher-advantage-only selectors to identify which criterion matters. A
Guided-OPD arm tests whether turn-level teacher anchoring already addresses the
failure. Match teacher, optimizer, action-token and teacher-scoring budgets,
environment interactions, and training-seed schedule; evaluate every arm from
the original unseeded state. The method claim fails if random or progress-
matched resets, TCOD-B2F, or Guided-OPD ties the targeted arm within the
preregistered smallest effect of interest.

\paragraph{Weights by recovery factorial.}
Evaluate base, round-2, and round-3 weights with recovery both off and on.
Preserve raw pre-rewrite emissions. This estimates main effects and interactions
for stage progress and format-defect rate instead of attributing a combined
system score to weights.

\paragraph{Generalization.}
Hold out at least one quest from candidate discovery, state selection, and
training, then compare full-walkthrough against no-walkthrough prompts. Report
the probability of reaching the bottleneck, crossing it conditional on arrival,
and finishing downstream. A second environment or model family is needed for
any broad claim. Corrected-interface SFT and a SCoRe-style first-error-prefix
condition are additional baselines under matched interaction and teacher-
scoring budgets \citep{lyu2025score}.

\paragraph{Copy-prior test.}
Construct paired canonical/malformed histories across tools, states, defect
types, and at least two teacher sizes or families. Score identical canonical
and malformed continuations; report teacher and student log-odds with paired
confidence intervals. Then intervene on grading context and test generation,
rather than inferring generation from grader-side movement.

\section{Related Work}

GKD establishes on-policy distillation from student-generated sequences with
teacher feedback \citep{agarwal2024gkd}; our work does not claim novelty for
OPD itself. TCOD directly studies temporal curricula and intermediate-state
initialization for multi-turn OPD \citep{wang2026tcod}; broad novelty for seeded
starts is therefore preempted. Guided-OPD changes live occupancy by mixing
teacher and student turns \citep{li2026guidedopd}; ReOPD designs reliable prefix
distributions offline \citep{liao2026reopd}; and SCoRe localizes the first
student error before short-horizon training \citep{lyu2025score}. SOD and KAT
instead alter or terminate supervision on degraded prefixes
\citep{zhong2026sod,xin2026kat}. These methods make intermediate-state,
student-failure, and state-centric novelty claims unavailable. DAgger formalizes
iterative data collection under the learner's state distribution
\citep{ross2011dagger}. Reset-along-demonstration, Backplay,
and Go-Explore show that restoring intermediate states can expose useful signal
in hard-exploration reinforcement learning
\citep{salimans2018montezuma,resnick2018backplay,ecoffet2021goexplore}. We use
the same control point for on-policy language-model distillation in a
persistent, typed environment, but the current study does not establish
superiority over those curriculum families.

Game-agent benchmarks already cover long-horizon language and vision agents.
BALROG emphasizes knowing--doing gaps across games \citep{paglieri2025balrog},
while Orak uses MCP across twelve games and releases gameplay fine-tuning data
\citep{park2026orak}. MCP and game play are therefore infrastructure here, not
novelty claims. Voyager and AgentTrek use curricula, reusable skills, or guided
replay to obtain long-horizon behavior
\citep{wang2023voyager,zheng2025agenttrek}. Our narrower question is whether a
frozen selector over low-occupancy, teacher-competent persistent external states
adds value beyond natural OPD, successful-prefix replay, turn-level guidance,
and matched generic resets.

Supervised imitation can compound errors under distribution shift
\citep{ross2011dagger}, and knowledge distillation does not guarantee functional
agreement between teacher and student \citep{stanton2021kd}. The r10
cross-vocabulary SFT regression in our development record is motivating
evidence, but its raw artifact bundle is incomplete and it is not a primary
result of this submission draft.

\section{Conclusion}

The observed case supports a concrete hypothesis: low-occupancy external states
may hide useful teacher competence from natural OPD. Natural-visitation OPD
changed behavior without changing score; a later hand-selected state-augmented
round was followed by unseeded passage of a previously unseen wall. That sequence
does not establish the proposed selector. Random and progress-matched resets,
TCOD-B2F, Guided-OPD, replication, and held-out transfer decide whether the
narrow method survives. A malformed-prefix probe separately shows why teacher-
forced token feedback can fail at the tool boundary.

\section*{Limitations}

The central limitation is statistical: one run per principal arm cannot support
a population claim, and three prompt variants are not independent seeds. The
state-seeded round also differs from the preceding round in initialization and
data. The environment is a single game with procedural walkthroughs, the
teacher and student are from one model family, and stage counts are coarse. Some
development analyses depend on logs whose raw inputs are not yet packaged. The
paper must not be submitted as confirmatory until the study design above is run
and a clean clone regenerates every main table.

Direct database construction is environment-specific and may not transfer to
systems without restorable state. The historical failure state was selected
retrospectively, so selection bias is possible; the prospective selector has no
outcome yet. Teacher competence at the target state is observed in a small number
of trajectories, not estimated precisely.
The syntax probe covers one defect family and should not be generalized to all
teacher-forced grading.

\section*{Ethical Considerations}

The work concerns game agents acting in isolated local worlds. Risks include
overstating autonomy, silently baking procedural knowledge into prompts, and
mistaking recovered actions for model competence. We mitigate these by labeling
scaffolding, logging model-visible schemas and raw emissions, separating runtime
recovery from weights, and reporting failed interventions. No human-subject or
private-user data are used.

\section*{Acknowledgments and AI Assistance Disclosure}

Large language models materially assisted literature search, code review,
experimental-design critique, and manuscript drafting. Human authors are
responsible for every claim, citation, artifact, and submission decision; the
scope will also be disclosed in the Responsible NLP Checklist.

\bibliography{submission}

\appendix

\section{Training Details}
\label{app:training}

The teacher is a scaffolded Qwen-family 4B model and the student begins from a
2B checkpoint. Both use the same vocabulary. Student rollouts are reconstructed
as the byte-exact serving prompt plus the verbatim emitted continuation. This
prefix property avoids re-rendering differences caused by chat-template handling
of reasoning delimiters. Teacher and behavior log probabilities are computed on
the identical raw string. Every tenth session is held out for development gates.

The update uses PPO-style clipped importance weighting with clip $\epsilon=0.3$,
advantage clamp $[-3,3]$, and a 1.5 weight on the first third of each session.
Training uses one epoch, learning rate $5\times10^{-5}$, effective batch size 32,
and a fresh LoRA initialized over the prior merged checkpoint. Action positions
alone receive the language-model head. These choices were development decisions,
not independently ablated.

Round 1 contains 5,564 train and 574 held-out records from the base policy's
natural rollouts. Round 2 combines 5,523 natural round-1 records and 2,326 seeded
records before the 7,024/825 train/held-out split. The seed encodes accepted
Herbalist stage 1, Foraging level 5, starter inventory, and a verified tile near
the target resource. Evaluation never loads this state.

\section{Claim-to-Evidence Boundary}

\begin{description}
\item[Round-1 behavior.] Descriptive, one run: measured behavior changed while score did not.
\item[State-seeded round.] Unreplicated: the seeded-training round was followed by unseeded wall passage in one run.
\item[Targeted selector.] Protocol only: no candidate-state bundle, trained checkpoint, or outcome is currently available.
\item[12 to 15.] Least-confounded historical weights-only observation; not an effect estimate, and exact parity remains unavailable.
\item[18/30.] Weights-plus-interface system result only.
\item[Copy prior.] Context-dependent endorsement reversal for the targeted historical probe.
\item[Recovery.] No repeat marker after correction in rewritten logs; raw relapse is unavailable.
\item[Cross-task transfer.] Unsupported; no claim.
\end{description}

\section{Reproducibility Contract}

Before confirmatory runs, each immutable run bundle must contain: environment and
game revisions; container or dependency locks; model, tokenizer, adapter, and
dataset revisions; complete model-visible tool schemas; rendered-prompt hashes;
sampling parameters; environment and inference seeds; raw emissions before any
rewrite; state-transition logs; analysis outputs; and checksums. Dataset records
must point to source runs and transformations. A clean-clone command must verify
hashes and regenerate every table and figure without network-dependent mutable
inputs.

The current development artifact does not yet meet this bar. In particular, raw
inputs for the historical r10 analysis and some OPD gate summaries are absent,
and the dedicated r10 eval path reset a different database from its game servers
after April 25. The June OPD runs used a separate database-aligned orchestrator
lane, but their raw bundles are also absent. Full dependency installation has
not been demonstrated from the repository root. These omissions are release
blockers, not appendix footnotes.

\section{Planned Statistical Analysis}

No session- or agent-level pseudo-replication will be used for policy effects.
The run is the analysis unit. A preliminary minimum of five runs per arm may be
used to estimate run-level variance, followed by a documented power calculation
for the confirmatory sample. The primary comparison is two-sided unless a
directional alternative and stopping rule are preregistered before collection.
We will report all runs, bootstrap or randomization-based confidence intervals
appropriate to the achieved sample, standardized and raw effect sizes, and the
full family of secondary endpoints with multiplicity handling. Exact categorical
wall-passage counts will be reported descriptively when sample sizes are small.

\section{Venue and Disclosure Note}

This PDF uses the current official ACL review style as an ARR/NAACL 2027 working
draft. The NAACL 2027 venue-specific call is not yet published. Recheck the call,
page limit, responsible NLP checklist, anonymity, author-review registration,
and commitment rules before submission. ARR prohibits simultaneous archival
review.

\end{document}
